\chapter{Credit and XVA}\label{ch:iv2:credit-and-xva}

\begin{ivtrains}
Hazard rates and credit default swaps by hand, then default correlation, then the exposure a derivatives desk has to its own counterparties and what that exposure costs. Every level here is invented for the arithmetic.
\end{ivtrains}

\begin{question}[$\star$ \iqroles{QR, QT, Trader, Risk} \iqfirm{bank, hedge fund} \iqlevel{junior}]\label{iq:iv2:credit-and-xva:1}
A five-year credit default swap on a company trades at 200 basis points. Assume recovery of 40\%. What default intensity, or hazard rate, does that imply? What is the chance that the company defaults within the five years?
\end{question}

\begin{question}[$\star$ \iqroles{QR, Risk, Analyst} \iqfirm{bank, hedge fund} \iqlevel{junior}]\label{iq:iv2:credit-and-xva:2}
A company has a constant hazard rate of 4\% a year. What is the chance that it defaults in the first year? What is the chance that it defaults in the second year, as seen from today? And in the second year, given that it survived the first?
\end{question}

\begin{question}[$\star$ \iqroles{QR, Trader, Analyst} \iqfirm{bank, hedge fund} \iqlevel{junior}]\label{iq:iv2:credit-and-xva:3}
The risk-free rate is 2\%, and a company has a hazard rate of 3\%. What is a five-year zero-coupon bond of face 1 worth with no recovery? What is it worth if 40\% of face is paid at maturity after a default? What yield spread does each one show?
\end{question}

\begin{question}[$\star$ \iqroles{Trader, Analyst, QT} \iqfirm{bank, hedge fund} \iqlevel{junior}]\label{iq:iv2:credit-and-xva:4}
You buy protection on \$10 million of a company at 200 basis points, paid quarterly. What do you pay each quarter? The company defaults and the auction sets a recovery of 30\%. What do you receive? Are you long or short credit, in the language of a credit desk?
\end{question}

\begin{question}[$\star$ \iqroles{QR, Trader, QT} \iqfirm{bank, hedge fund} \iqlevel{junior}]\label{iq:iv2:credit-and-xva:5}
A five-year CDS pays a fixed coupon of 100 basis points, and the name trades at a spread of 200. The rate is 3\% and the recovery 40\%. Who pays whom up front, and about how many points? Take a continuous premium and a flat hazard rate.
\end{question}

\begin{question}[$\star\star$ \iqroles{QR, Trader, Risk} \iqfirm{bank, hedge fund} \iqlevel{junior, experienced}]\label{iq:iv2:credit-and-xva:6}
With a fixed coupon of 100 basis points, a rate of 3\% and five years, work out the up-front payment for a name at 200 and for a name at 2\,000, once at a recovery of 40\% and once at 20\%. Why can you ignore the recovery assumption for the first name but not for the second?
\end{question}

\begin{question}[$\star\star$ \iqroles{Trader, QT, QR} \iqfirm{bank, hedge fund} \iqlevel{junior, experienced}]\label{iq:iv2:credit-and-xva:7}
A floating-rate bond of a company pays a spread of 220 basis points over the funding rate and trades at par. The five-year CDS on the same company trades at 180. What trade do you put on, and what does it earn on \$10 million? Why do the two spreads have to match in a clean model? What can go wrong with the trade?
\end{question}

\begin{question}[$\star\star$ \iqroles{QR, Risk} \iqfirm{bank, hedge fund} \iqlevel{junior, experienced}]\label{iq:iv2:credit-and-xva:8}
Two companies default within a year with probabilities 5\% and 5\%, and their default correlation is 0.3. What is the chance that both default? That at least one does? Now take probabilities of 5\% and 1\%. How high can the default correlation possibly go?
\end{question}

\begin{question}[$\star\star$ \iqroles{QR, Trader, Risk} \iqfirm{bank, hedge fund} \iqlevel{junior, experienced}]\label{iq:iv2:credit-and-xva:9}
A basket holds five names, each with a one-year default probability of 1\% and recovery of 40\%. A first-to-default swap pays the loss on the first name to default, and then stops. About what spread is it worth if the defaults are independent? If they are perfectly correlated? What happens to a second-to-default as correlation rises?
\end{question}

\begin{question}[$\star\star$ \iqroles{QR, Risk} \iqfirm{bank} \iqlevel{junior, experienced}]\label{iq:iv2:credit-and-xva:10}
You have a five-year at-the-money interest rate swap on a notional of 100, with no discounting. The par rate moves as a driftless Brownian motion with a volatility of 1\% a year. Write the expected positive exposure at time $t$. When is it largest, and how large is it?
\end{question}

\begin{question}[$\star\star$ \iqroles{QR, Risk} \iqfirm{bank} \iqlevel{junior, experienced}]\label{iq:iv2:credit-and-xva:11}
A five-year swap has expected exposure $0.399 \times 0.01 \times 100 \sqrt{t}\,(5 - t)$ on a notional of 100, and the counterparty has a hazard rate of 2\% and recovery of 40\%. Ignoring discounting and the small decay in survival, what is the credit valuation adjustment, or CVA? Show that it is about the counterparty's spread times the average exposure times the life.
\end{question}

\begin{question}[$\star\star$ \iqroles{QR, Risk, QD} \iqfirm{bank} \iqlevel{junior, experienced}]\label{iq:iv2:credit-and-xva:12}
Two trades face the same counterparty. At a given date, each has a normal mark to market with mean 0 and standard deviation 1. What is the expected positive exposure with no netting? What is it with netting, at correlations of $-1$, $-0.5$, 0, $0.5$ and 1?
\end{question}

\begin{question}[$\star\star$ \iqroles{QR, Risk} \iqfirm{bank} \iqlevel{junior, experienced}]\label{iq:iv2:credit-and-xva:13}
A trade's value moves with a normal volatility of \$10 million a year, and has mean 0. What is its expected exposure at one year with no collateral? It now has daily variation margin with zero threshold, and it would take 10 business days to close out after a default. What does the exposure become? Why is it not zero?
\end{question}

\begin{question}[$\star\star\star$ \iqroles{QR, Risk} \iqfirm{bank} \iqlevel{experienced}]\label{iq:iv2:credit-and-xva:14}
A five-year at-the-money swap on 100 has a symmetric expected exposure, positive and negative, each with a profile that integrates to 5.947 over the life. The counterparty's hazard is 2\%, yours is 1\%, and both recover 40\%. Work out CVA and the debit valuation adjustment, or DVA, without and with first-to-default weighting. What happens to your profit when your own credit gets worse? Why do many desks refuse to treat that as real money?
\end{question}

\begin{question}[$\star\star\star$ \iqroles{QR, Risk, Trader} \iqfirm{bank} \iqlevel{experienced}]\label{iq:iv2:credit-and-xva:15}
You trade an uncollateralised five-year swap with a client and hedge it with a fully collateralised one, on 100 of notional, with an integrated exposure of 5.947. You borrow at 100 basis points over the collateral rate. What is the funding valuation adjustment, or FVA? Now suppose your whole funding spread is credit: 60 basis points, from a hazard of 1\% and recovery of 40\%. Compare the funding benefit on negative exposure with your DVA. What do you conclude?
\end{question}

\begin{question}[$\star\star\star$ \iqroles{QR, Risk, Trader} \iqfirm{bank} \iqlevel{experienced}]\label{iq:iv2:credit-and-xva:16}
Your exposure to a counterparty in one year is $\max(-5Z, 0)$, with $Z$ standard normal. The counterparty's one-year default probability is 2\%, and it defaults when $\rho Z + \sqrt{1 - \rho^2}\,\varepsilon$ falls below the threshold that gives that probability. Find the expected exposure given default, at $\rho$ of 0, 0.3, 0.5 and $-0.3$. Give a real example of each sign.
\end{question}

\begin{question}[$\star\star\star$ \iqroles{QR, Trader} \iqfirm{bank, hedge fund} \iqlevel{experienced}]\label{iq:iv2:credit-and-xva:17}
A large pool of identical names has a default probability of 2\% and recovery of 40\%, under a one-factor Gaussian copula with correlation $\rho$. Give the expected loss of the tranches 0--3\%, 3--7\%, 7--15\% and 15--100\% as a share of each tranche, at $\rho$ of 0.05, 0.2 and 0.5. Which way does each move? What stays fixed?
\end{question}

\begin{question}[$\star\star\star$ \iqroles{QR, Analyst, Risk} \iqfirm{bank, hedge fund} \iqlevel{junior, experienced}]\label{iq:iv2:credit-and-xva:18}
A firm has assets of 100 with a volatility of 25\%, and zero-coupon debt of face 80 due in one year. The rate is 3\%. In the Merton model, what are the equity value, the debt value, the risk-neutral default probability and the credit spread? What does the spread do at three months and at five years? What does that tell you about the model?
\end{question}

\begin{question}[$\star\star\star$ \iqroles{QR, QD, Risk} \iqfirm{bank, hedge fund} \iqlevel{experienced}]\label{iq:iv2:credit-and-xva:19}
CDS spreads at one, three and five years are 100, 150 and 200 basis points, the rate is 3\% and the recovery 40\%. Find the hazard rates for the three pieces, 0--1, 1--3 and 3--5, first from the credit triangle and then by a proper bootstrap with quarterly premiums. Why is the last piece higher by the proper route? What would a spread curve of 2\,000 at one year and 800 at five years be telling you?
\end{question}

\begin{question}[$\star\star\star\star$ \iqroles{QR, Trader, Risk} \iqfirm{bank} \iqlevel{experienced}]\label{iq:iv2:credit-and-xva:20}
You run the CVA desk. You hold an uncollateralised five-year at-the-money swap on \$100 million with a corporate, rate volatility 1\% a year, whose five-year CDS is at 200 basis points with recovery of 40\%. There is no discounting. Value the CVA with survival weighting. How much CDS protection do you buy to hedge the spread risk? Now the swap moves 50 basis points in your favour, and separately the spread goes to 300. What is the CVA after each, and after both? What does the gap between both and the sum of the two moves mean for your hedge? What would you do if the corporate has no traded CDS?
\end{question}
